\documentclass[10.5pt]{article}
\usepackage[letterpaper, margin=0.6in, top=0.5in, bottom=0.5in]{geometry}
\usepackage{amsmath,amssymb,amsthm}
\usepackage{fancyhdr}
\usepackage{xcolor}
\usepackage{tikz}
\usepackage{enumitem}
\usepackage{tcolorbox}
\usepackage{hyperref}

% Define colors matching instructor board annotations
\definecolor{chalkblue}{RGB}{0, 105, 185}
\definecolor{chalkred}{RGB}{205, 35, 35}
\definecolor{chalkpink}{RGB}{230, 90, 120}
\definecolor{chalkcyan}{RGB}{15, 160, 190}
\definecolor{chalkgreen}{RGB}{20, 130, 50}
\definecolor{boxbg}{RGB}{252, 252, 254}
\definecolor{boxframe}{RGB}{150, 160, 180}

\pagestyle{fancy}
\fancyhf{}
\renewcommand{\headrulewidth}{0pt}
\renewcommand{\footrulewidth}{0pt}

\setlength{\parindent}{0pt}
\setlength{\parskip}{2.5pt}

\begin{document}

% ==============================================================================
% PAGE 1
% ==============================================================================
\noindent \textbf{\large Math 105L} \hfill \textbf{\large Lecture 1-3}\\[1pt]
{\textbf{\Large Linear Functions: Slope Characterization \& Applied Modeling}}\\[-5pt]
\rule{\textwidth}{1.2pt}

\vspace{0.03in}

\textbf{\large 1. Fundamental Characterization of Linear Slope}

The slope $m$ of a linear relationship connecting variables $x$ (independent) and $y$ (dependent) is given by:
\[
  m = \frac{\Delta y}{\Delta x} = \frac{y_2 - y_1}{x_2 - x_1}
\]
Setting $\mathbf{\Delta x = 1}$ yields:
\[
  m = \frac{\Delta y}{1} \implies \mathbf{\Delta y = m}
\]

\begin{tcolorbox}[colback=boxbg, colframe=boxframe, boxrule=0.8pt, arc=3pt, left=8pt, right=8pt, top=4pt, bottom=4pt]
\textbf{The Universal Operational Meaning of Slope:}\\[2pt]
The slope $m$ represents the exact amount the \textbf{dependent variable ($y$)} changes every time the \textbf{independent variable ($x$)} increases by \textbf{one unit}.
\begin{itemize}[leftmargin=1.2em, topsep=1pt, itemsep=2pt]
  \item If $m > 0$, the dependent variable increases by $m$ per unit increase in input.
  \item If $m < 0$, the dependent variable decreases by $|m|$ per unit increase in input.
  \item \textbf{Units of Slope:} Always $\mathbf{\frac{\text{Units of } y}{\text{Units of } x}}$.
\end{itemize}
\end{tcolorbox}

\vspace{0.04in}

\textbf{\large 2. Classroom Case Study: Hanford Cancer Mortality Modeling}

From the Cancer Mortality laboratory, linear regression connects the radioactive exposure index $I$ to annual cancer mortality rate $M$ (measured in cancer deaths per 20,000 residents per year):
\[
  \mathbf{M = 11.2 I + 112}
\]
\textbf{Slope Interpretation:} For each 1-unit increase in the exposure index $I$, cancer mortality increases by $11.2$ deaths per 20,000 residents each year.

\vspace{0.03in}

\begin{tcolorbox}[colback=chalkblue!6, colframe=chalkblue, boxrule=1pt, arc=3pt, left=8pt, right=8pt, top=4pt, bottom=4pt]
\textbf{Classroom Problem Worked in Full:}\\[2pt]
A mayor of a town with population $P = 115,000$ enacts remediation policies that successfully decrease the local exposure index by $\mathbf{\Delta I = -1.7}$ units. How many cancer deaths does the model predict will be prevented over the next $\mathbf{5}$ years?
\end{tcolorbox}

\begin{itemize}[leftmargin=1.2em, topsep=1pt, itemsep=4pt]
  \item \textbf{Step 1: Calculate the reduction in mortality rate per 20,000 people:}
  \[
    \Delta M = m \cdot \Delta I = 11.2 \cdot (-1.7) = \mathbf{-19.04 \text{ deaths per 20,000 people per year}}
  \]
  \item \textbf{Step 2: Scale to the town's total population ($P = 115,000$):}
  Determine the number of 20,000-person population blocks:
  \[
    \text{Blocks} = \frac{115,000}{20,000} = \mathbf{5.75 \text{ blocks}}
  \]
  \item \textbf{Step 3: Compute annual deaths prevented town-wide:}
  \[
    \text{Deaths prevented per year} = 19.04 \times 5.75 = \mathbf{109.48 \text{ deaths/year}}
  \]
  \item \textbf{Step 4: Compute cumulative deaths prevented over a 5-year span:}
  \[
    \text{Total deaths prevented} = 109.48 \text{ deaths/year} \times 5 \text{ years} = \mathbf{547.4 \approx 547 \text{ lives saved}} \qquad \checkmark
  \]
\end{itemize}

\end{document}

